% monstyle.tex - Style & Fonctions pour TikZ
\usepackage{tikz}
\usepackage{xcolor}
\usepackage{amsmath}
\usepackage{anyfontsize}
\usepackage{bm}
\usetikzlibrary{3d, perspective}
\usetikzlibrary{arrows.meta, calc, positioning}
\usetikzlibrary{hobby} % <-- indispensable pour use Hobby shortcut

% --- Palette de Couleurs (Style Mer de Fermi) ---
\definecolor{colorOne}{HTML}{443E46}   % Anthracite principal (texte, axes)
\definecolor{colorTwo}{HTML}{FFFFFF}   % Blanc
\definecolor{colorThree}{HTML}{908CA4} % Gris-bleu (grille, repères)
\definecolor{colorFour}{HTML}{57659E}  % Bleu indigo (courbe 1, deltas)
\definecolor{colorFive}{HTML}{C57284}  % Rose/Magenta doux
\definecolor{colorSix}{HTML}{FF5B69}   % Rouge/Corail vif (courbe 2, points)
\definecolor{colorGold}{HTML}{FFD700}

% Raccourcis pour les couleurs
\def\colorOne{colorOne}
\def\colorTwo{colorTwo}
\def\colorThree{colorThree}
\def\colorFour{colorFour}
\def\colorFive{colorFive}
\def\colorSix{colorSix}
\def\colorGold{colorGold}

% Définition des couleurs avec les codes HTML
\definecolor{colorOne}{HTML}{443E46}
%\definecolor{colorTwo}{HTML}{F6DEB8}
\definecolor{colorTwo}{HTML}{FFFFFF}

\definecolor{colorThree}{HTML}{908CA4}
\definecolor{colorFour}{HTML}{57659E}
\definecolor{colorFive}{HTML}{C57284}
\definecolor{colorSix}{HTML}{FF5B69}

% Raccourcis pour les couleurs
\def\colorOne{colorOne}
\def\colorTwo{colorTwo}
\def\colorThree{colorThree}
\def\colorFour{colorFour}
\def\colorFive{colorFive}
\def\colorSix{colorSix}

% --- Style Global TikZ ---
\tikzset{
    every picture/.style={
        color=colorOne,
        font=\fontsize{12pt}{14pt}\selectfont\bfseries
    },
    fermi axis/.style={
        thick,
        line width=0.6ex,
        ->,
        >=Stealth,
        color=colorOne
    },
    wave primary/.style={
        line width=0.8ex,
        color=colorFour
    },
    wave secondary/.style={
        line width=0.8ex,
        color=colorSix
    }
}

% --- Commande de Grille Personnalisée ---
% Usage : \drawgrid{xmin}{xmax}{ymin}{ymax}[pas]
% \newcommand{\drawgrid}[4][1]{
%     \begin{scope}[opacity=0.25]
%         \draw[very thin, color=colorThree, step=0.2*#1] (#1*#1 - #1, #3) grid (#2, #4);
%         \draw[thin, color=colorThree, step=#1] (#1*#1 - #1, #3) grid (#2, #4);
%     \end{scope}
% }

% --- Commande d'Axes Typés ---
\newcommand{\drawCustomAxes}[4]{
    % #1=xmin, #2=xmax, #3=ymin, #4=ymax
    \draw[fermi axis] (#1, 0) -- (#2, 0) node[pos=1.01, below, color=colorOne] {\huge $x$};
    \draw[fermi axis] (0, #3) -- (0, #4) node[pos=1.01, left, color=colorOne] {\huge $y$};
}

% Définition de la fonction pour générer la grille
% #1 = xmin
% #2 = xmax
% #3 = ymin
% #4 = ymax
% #5 = échelle X
% #6 = échelle Y
\newcommand{\drawgrid}[6]{%

    \begin{scope}[
        opacity=0.2,
        xscale=#5,
        yscale=#6
    ]

        % Grille principale
        \draw[
            very thin,
            gray
        ]
        (#1,#3)
        grid
        (#2,#4);

        % Sous-grille
        \draw[
            very thin,
            gray
        ]
        (#1,#3)
        grid[
            step=0.1
        ]
        (#2,#4);

        % Axe X
        \draw[
            thick,
            line width=0.8ex,
            ->,
            >=Stealth,
            black
        ]
        (#1,#3)
        --
        (#2,#3);

        % Graduations X
        \foreach \x in {#1,...,#2} {
            \draw[
                line width=0.8ex
            ]
            (\x,#3-0.1)
            --
            (\x,#3+0.1);

            \node[
                below
            ]
            at (\x,#3)
            {\x};
        }

        % Axe Y
        \draw[
            thick,
            line width=0.8ex,
            ->,
            >=Stealth,
            black
        ]
        (#1,#3)
        --
        (#1,#4);

        % Graduations Y
        \foreach \y in {#3,...,#4} {
            \draw[
                line width=0.8ex
            ]
            (#1-0.1,\y)
            --
            (#1+0.1,\y);

            \node[
                left
            ]
            at (#1,\y)
            {\y};
        }

    \end{scope}
}


\newcommand{\Axex}[1][$x$]{ % "x" est la valeur par défaut de #1
    \draw
        (-10, 0) edge[thick, line width=0.8ex, ->, >=Triangle, color=\colorThree] 
        node[shift = {(0,1)} , pos=1.01, below]{ \huge #1 } (10, 0)
        ;
}

\newcommand{\Axetheta}[1][$\theta$]{
	\draw
		(-10, 0 ) edge [thick,line width=0.8ex,->,>=Triangle , color = \colorThree ]node [pos=1.01,right]{\huge #1 } ( 10  , 0 )
	;
}

\newcommand{\Axedensity}[1][$n$]{
	\draw
		(0, -1 ) edge [thick,line width=0.8ex,->,>=Triangle , color = \colorThree ]node [pos=1.01,left  ]{\huge #1 } ( 0  , 6 )
	;
}

\newcommand{\Axedistribution}[1][$\nu$]{
	\draw
		(0, -1 ) edge [thick,line width=0.8ex,->,>=Triangle , color = \colorThree ]node [pos=1.01,left  ]{\huge #1 } ( 0  , 6 )
	;
}

\newcommand{\Axesphase}[2][$x$]{
	\Axex[#1]
	\begin{scope}[rotate = 90]
		\Axetheta[#2]	
	\end{scope}
}

\newcommand{\Axesdensity}[2][$x$]{
	\Axex[#1]
	\Axedensity[#2]
}

\newcommand{\Axesdistribution}[2][\theta]{
	\Axetheta[#1]
	\Axedistribution[#2]
}

\newcommand{\Axetemps}[1][Deformation]{
	\draw 
		(-5 , 0 ) edge [line width=2.8ex, color=\colorOne ]( 0, 0)
		(43 , 0 ) edge [line width=2.8ex, color=\colorOne , ->, >=Triangle] node[pos = 1.0 , right , scale = 2 ] { \fontsize{60pt}{720pt}\selectfont $\tilde{t}$} ( 50 , 0)
		( 0 , 1 ) edge [line width=2.0ex, color=\colorOne ] ( 0 , -1 )  
		( 43 , 1 ) edge [line width=2.0ex, color=\colorOne ] ( 43 , -1 ) 
		(0 , 0) edge[<->, > = stealth, line width=2.8ex, color=\colorOne ]  node [ rectangle , midway, fill = white , scale = 2 ] {\fontsize{600pt}{720pt}\selectfont \bf  #1}(43, 0)  
	;

}

\newcommand{\Palette}[1][\colorOne]{
	\clip[decorate, decoration={random steps, segment length=3pt, amplitude=3pt}]
        	(-1,-1) rectangle (1,1);
	\draw[fill=#1 , color = #1 ] (-2,-2) rectangle (2,2);
}

% Graduation ordonné

\newcommand{\GraduationYMm}[2]{
	\foreach \r in {1,...,9} {
		\draw 
			({-0.25 * 0.5 * (cos(3.14*0.2*\r r)^2 + 1)}, #1*#2 - #2 + 0.1*#2*\r) 
			edge [thick, line width=0.3ex, color=\colorOne] 
			({0.25 * 0.5 * (cos(3.14*0.2*\r r)^2 + 1)}, #1*#2 - #2 + 0.1*#2*\r);
	}
}

\newcommand{\GraduationYCm}[4]{% #1: min, #2: max, #3: liste
  \foreach \R in {#1,...,#2} {
    \pgfmathparse{#3[\R-(#1)]} % index dans la liste
    \let\labelText\pgfmathresult

    \draw[color=\colorThree]
      (-0.3, \R*#4)
      edge [thick, line width=0.5ex, color=\colorOne]
      node [pos=-0.0, left]{\Large \labelText}
      (0.3, \R*#4);

    %\GraduationYMm{\R}{#4}
  }

  % Graduation supplémentaire à #2 + 3
  \pgfmathparse{#2 + 1}
  \edef\extendedPosition{\pgfmathresult}
  %\GraduationYMm{\extendedPosition}{#4}
}



\newcommand{\GraduationXMm}[2]{
	\foreach \r in {1,...,9} {
		\draw 
			( #1*#2 - #2 + 0.1*#2*\r,{-0.25 * 0.5 * (cos(3.14*0.2*\r r)^2 + 1)}) 
			edge [thick, line width=0.3ex, color=\colorOne] 
			( #1*#2 - #2 + 0.1*#2*\r,{0.25 * 0.5 * (cos(3.14*0.2*\r r)^2 + 1)});
	}
}

\newcommand{\GraduationXCm}[4]{% #1: min, #2: max, #3: liste
  \foreach \R in {#1,...,#2} {
    \pgfmathparse{#3[\R-(#1)]} % index dans la liste
    \let\labelText\pgfmathresult

    \draw[color=\colorThree]
      (\R*#4,-0.3)
      edge [thick, line width=0.5ex, color=\colorOne]
      node [pos=-0.0, below]{\Large \labelText}
      (\R*#4,0.3);

    %\GraduationXMm{\R}{#4}
  }

  % Graduation supplémentaire à #2 + 3
  \pgfmathparse{#2 + 1}
  \edef\extendedPosition{\pgfmathresult}
  %\GraduationXMm{\extendedPosition}{#4}
}

\newcommand{\GraduationsY}[6]{%
    % Force l'expansion de la liste
    \edef\graduationsList{#1}%

    \foreach \position/\labelLeft/\labelRight in \graduationsList {%

        % Calcul de la position réelle
        \pgfmathsetmacro{\posY}{\position * #2}%

        % Trait de graduation
        \draw[
            thick,
            line width=#4,
            color=#5
        ]
        (-#3,\posY)
        --
        (#3,\posY);

        % Label gauche
        \ifx\labelLeft\empty
        \else
            \node[
                left
            ] at (-#3*1.5,\posY)
            {#6 \labelLeft};
        \fi

        % Label droit
        \ifx\labelRight\empty
        \else
            \node[
                right
            ] at (#3*1.5,\posY)
            {#6 \labelRight};
        \fi
    }%
}

\newcommand{\GraduationsX}[6]{%
    % Force l'expansion de la liste
    \edef\graduationsList{#1}%

    \foreach \position/\labelBottom/\labelTop in \graduationsList {%

        % Calcul de la position réelle
        \pgfmathsetmacro{\posX}{\position * #2}%

        % Trait de graduation
        \draw[
            thick,
            line width=#4,
            color=#5
        ]
        (\posX,-#3)
        --
        (\posX,#3);

        % Label bas
        \ifx\labelBottom\empty
        \else
            \node[
                below
            ] at (\posX,-#3*1.5)
            {#6 \labelBottom};
        \fi

        % Label haut
        \ifx\labelTop\empty
        \else
            \node[
                above
            ] at (\posX,#3*1.5)
            {#6 \labelTop};
        \fi
    }%
}

% Définition de la fonction pour générer la grille
\newcommand{\backgrowngrid}[9]{


	\begin{scope}[xscale = #5 , yscale = #6 ]


    	% Paramètres : #1=xmin, #3=xmax, #2=ymin, #4=ymax
    	%\clip[scale = 0.95 , decorate, decoration={random steps, segment length=0.1 cm, amplitude=0.2cm}](#1 , #2 ) rectangle ( #3 , #4 ) ;
    	\pgfmathparse{#1-1}  
    	\let\valXmin\pgfmathresult
    	\pgfmathparse{#2-1}  
    	\let\valYmin\pgfmathresult
    	\pgfmathparse{#3+1}  
    	\let\valXmax\pgfmathresult
    	\pgfmathparse{#4+1}  
    	\let\valYmax\pgfmathresult

    	\draw[fill = #7, opacity = 0.1]  (\valXmin , \valYmin ) rectangle ( \valXmax , \valYmax ) ;

   	 	% Grille principale en cm
    	\draw[very thin, line width = #9 ,  color = #8] (\valXmin , \valYmin ) grid ( \valXmax , \valYmax ); % Grille de base (#5 cm , #6 cm )


    \end{scope}


}



% ============================================================
% Figures TikZ
% ============================================================

\usepackage{graphicx}
\usepackage{caption}
\usepackage{etoolbox}


% ------------------------------------------------------------
% Contenu d'une figure TikZ
% ------------------------------------------------------------
% #1 = ref
% #2 = scale
%
% Si scale est vide :
%     \input{figures/tikz/xxx.tex}
%
% Sinon :
%     \scalebox{scale}{\input{...}}

\newcommand{\InsertTikzFigureContent}[2]{%
    \ifstrempty{#2}{%
        \input{figures/tikz/#1.tex}%
    }{%
        \scalebox{#2}{%
            \input{figures/tikz/#1.tex}%
        }%
    }%
}


% ------------------------------------------------------------
% Figure normale
% ------------------------------------------------------------
% #1 = ref
% #2 = scale
% #3 = légende
% #4 = label

\newcommand{\InsertTikzFigureNormal}[4]{%
    \begin{figure}[h!]
        \centering

        \InsertTikzFigureContent{#1}{#2}

        \ifstrempty{#3}{}{%
            \caption{#3}%
        }

        \ifstrempty{#4}{}{%
            \label{#4}%
        }

    \end{figure}%
}


% ------------------------------------------------------------
% Figure dans un encadré
% ------------------------------------------------------------
% #1 = ref
% #2 = scale
% #3 = légende
% #4 = label

\newcommand{\InsertTikzFigureInline}[4]{%
    \begin{center}
        \nopagebreak

        \InsertTikzFigureContent{#1}{#2}

        \ifstrempty{#3}{}{%
            \captionof{figure}{#3}%
        }

        \ifstrempty{#4}{}{%
            \label{#4}%
        }

    \end{center}%
}